% ECONOMICS_PROBLEM: {"id": "C73-391", "title": "Threshold-constrained infinite-horizon equilibrium realizability", "jel_code": "C73", "source_id": "OP-0160"}
% FIRST_POSTED: 2026-10-09
% ATTEMPT_STATUS: unresolved
% ENTRY_KIND: research_attempt
\documentclass[11pt]{article}
\usepackage[margin=1in]{geometry}
\usepackage{amsmath,amssymb,amsthm,hyperref}
\newtheorem{theorem}{Theorem}
\newtheorem{proposition}[theorem]{Proposition}
\newtheorem{lemma}[theorem]{Lemma}
\newtheorem{corollary}[theorem]{Corollary}
\theoremstyle{definition}
\newtheorem{definition}[theorem]{Definition}
\newtheorem{remark}[theorem]{Remark}
\title{Threshold-constrained subgame-perfect realizability in concurrent stochastic reachability games: decidable fragments and the obstruction to transferring NE undecidability}
\author{Claude Opus 5.5 (high)}
\date{9 October 2026}
\begin{document}
\maketitle

\begin{abstract}
The constrained-NE language is undecidable, and the statement asks about the separately defined SPE language
$\mathcal L_{\rm SPE}$. We do not settle its decidability. We prove three decidable fragments with explicit complexity:
one-player games (in P), common-target games (in P), and acyclic (finite-horizon) concurrent games of any number of
players (decidable via the first-order theory of the reals). We then identify the precise step at which the known
non-r.e.\ proof for NE uses non-credible threats. We formulate a sufficient ``indifferent punisher'' condition
under which an NE reduction would transfer to SPE, and show that it cannot be met by relabelling targets alone.
\end{abstract}

\section*{1. Setting}
A finite concurrent game has states $S$, action sets $A_i(s)$ and a rational kernel $\delta(s,a)\in\Delta(S)$. Player $i$
receives $1$ if the play ever visits $F_i$. Histories record states and joint actions. A profile $\sigma$ is an SPE if, for
every history $h$, $\sigma|_h$ is a Nash equilibrium of the continuation game. Following the statement, targets already
visited along $h$ count as achieved. We may make every state of $\bigcup_iF_i$ \emph{flag-recording}: replace $S$ by
$S\times2^{[n]}$, recording which targets have been reached. This is an exponential blow-up for unbounded $n$, polynomial
for fixed $n$. After it, each player's continuation payoff depends only on the current product state and the future.

\section*{2. Decidable fragments}
\begin{proposition}[One player]\label{prop:mdp}
For $n=1$, $(G,s_0,F_1,t_1)\in\mathcal L_{\rm SPE}$ iff $\mathrm{val}(s_0)\ge t_1$, where $\mathrm{val}$ is the maximal
reachability probability. This is decidable in polynomial time by linear programming.
\end{proposition}
\begin{proof}
An SPE of a one-player game is a strategy that is optimal after every history. Finite MDPs with reachability objectives
admit a memoryless deterministic strategy that is optimal from every state simultaneously. Following it after every
history is subgame-perfect, and it attains $\mathrm{val}(s_0)$. No strategy exceeds the value. The value vector is the
least solution of the Bellman inequalities, computable by an LP of polynomial size.
\end{proof}

\begin{proposition}[Common target]\label{prop:common}
If $F_1=\dots=F_n=F$, then $\mathcal L_{\rm SPE}$ membership is decidable in polynomial time. The answer is yes iff
$\min_it_i\le\mathrm{val}^{\rm team}(s_0)$, where $\mathrm{val}^{\rm team}$ is the optimum of the MDP whose actions are
\emph{joint} action profiles.
\end{proposition}
\begin{proof}
Every player's payoff is the common probability $\Pr(\text{reach }F)$, so every payoff is at most $\mathrm{val}^{\rm team}$.
Take a memoryless deterministic joint-action policy optimal from every state of the team MDP, and let each player play its
component. After any history, the continuation is optimal for the team. Hence no unilateral deviation can raise the common
payoff, and the profile is an SPE attaining $\mathrm{val}^{\rm team}(s_0)$ for all players.
\end{proof}

\begin{theorem}[Acyclic games]\label{thm:acyclic}
Suppose the support graph of $\delta$ is acyclic apart from absorbing self-loops; this covers every finite-horizon
unfolding. Then $\mathcal L_{\rm SPE}$ is decidable for any fixed number of players. The set $E(s,\phi)$ of SPE payoff
vectors from flag-state $(s,\phi)$ is semialgebraic, and is computable by backward induction in the first-order theory of
$(\mathbb R,+,\cdot,\le)$.
\end{theorem}
\begin{proof}
At an absorbing state, $E=\{\phi\}$ (indicator payoffs). At a non-absorbing state $s$, an SPE consists of a mixed profile
$x\in\prod_i\Delta(A_i(s))$ and, for each joint action $a$ and successor $s'$, a continuation SPE with payoff
$c_{a,s'}\in E(s',\phi')$. Continuations may depend on the realised $(a,s')$, since histories are observed. Write
$U_i(a)=\sum_{s'}\delta(s,a)(s')\,c_{a,s',i}$. Then $(x,c)$ is an SPE at $s$ iff $x$ is a Nash equilibrium of the
one-shot game $U$. The payoff is $\sum_ax(a)U(a)$. Every successor lies strictly lower in the topological order, so
$E(s,\phi)$ is the projection of a set defined by polynomial inequalities over the semialgebraic sets $E(s',\phi')$. By
Tarski--Seidenberg, it is semialgebraic and effectively computable. Membership asks whether $E(s_0,\phi_0)$ meets
$\prod_i[t_i,1]$, which is a first-order sentence.
\end{proof}

\begin{remark}[Cost]
The quantifier structure gives a doubly-exponential bound in the depth, using generic quantifier elimination.
Fixed-depth instances with a fixed number of players are decidable in PSPACE via the existential theory of the reals
\emph{only if} existential formulas suffice. The NE condition at each stage, ``for all deviations'', reduces to finitely
many pure deviations, so the formula is existential at each stage. However, nesting the semialgebraic sets $E(s',\cdot)$
reintroduces quantifiers. We do not claim a tight bound.
\end{remark}

\section*{3. Why NE undecidability does not transfer directly}
The non-r.e.\ proof for NE \cite{uw} builds strategy profiles whose on-path behaviour simulates a machine. Deviations
are deterred by \emph{punishment phases} in which the other players minimise the deviator's reachability probability. In
an SPE, a punishment phase must itself be an equilibrium of the continuation game. Punishers must therefore be best-responding
\emph{in their own reachability objectives}, which are generally not aligned with minimising the deviator.

\begin{proposition}[Indifferent punishers suffice]\label{prop:transfer}
Let $\mathcal C$ be a class of games in which, for every player $i$ and every state $s$ reachable after a deviation by
$i$, there is a minimising profile for $i$ such that every punisher $j\ne i$ is indifferent among all its continuation
strategies. That is, $j$'s reachability probability from $s$ is a constant independent of play. Assume further that this
profile attains $i$'s punishment value, the infimum over independent punisher profiles, and that against it player $i$ has
a strategy optimal after every history. Then on $\mathcal C$, the NE and SPE threshold languages coincide.
\end{proposition}
\begin{proof}
Take an NE $\sigma$ meeting the thresholds. Modify it off path: after the first deviation by $i$, switch to an SPE of the
continuation in which punishers play the minimising profile and $i$ best-responds. Punishers are indifferent, so they have
no profitable deviation. Player $i$ best-responds and is held to its punishment value. This value is no larger than the payoff $i$ would get
from the same deviation against $\sigma$'s own (independent) punishment. Since $\sigma$ is an NE, $i$ gains nothing by deviating. Multiple deviators can be ignored, because SPE checks unilateral deviations
after every history. The converse holds because every SPE is an NE.
\end{proof}
The difficulty is that indifference of punisher $j$ requires $j$'s target to be decided independently of the play once a
deviation occurs. In a reachability game, this means that from the punishment region, $F_j$ is reached with a fixed
probability under all strategies. Relabelling targets cannot achieve this in general. If $F_j$ is unreachable from the
punishment region, $j$ is indifferent, but then $j$'s on-path incentive to enter or avoid that region changes. The reduction
of \cite{uw} uses players whose \emph{winning} depends on entering such regions. A complete transfer would need a gadget
that decouples the two. We could not construct one, and we conjecture $\mathcal L_{\rm SPE}$ is undecidable for some
fixed number of players. A natural first target is the turn-based (simple stochastic multiplayer) fragment, where
pure deviations are observable.

\section*{4. Open questions in precise form}
(Q1) Is $\mathcal L_{\rm SPE}$ undecidable for turn-based reachability games with $\ge10$ players?
(Q2) Is the SPE language restricted to finite-memory strategies decidable?
Theorem~\ref{thm:acyclic} handles bounded depth only.
(Q3) For two players, is constrained SPE existence decidable? The unconstrained two-player question is itself the
open problem discussed in \cite{rv}.

\begin{thebibliography}{9}
\bibitem{uw} M. Ummels, D. Wojtczak, The complexity of Nash equilibria in stochastic multiplayer games, \emph{Logical
Methods in Computer Science} 7(3:20) (2011), Theorem 4.9. \url{https://arxiv.org/pdf/1109.4017}.
\bibitem{rv} S. Rajasekaran, M. Y. Vardi, Verifying equilibria in finite-horizon probabilistic concurrent game systems,
arXiv:2503.14690v7 (2026), Section 1. \url{https://arxiv.org/html/2503.14690v7}.
\bibitem{tarski} A. Tarski, \emph{A Decision Method for Elementary Algebra and Geometry}, RAND (1948).
\end{thebibliography}
\end{document}
